\documentclass[conference]{IEEEtran}

\usepackage{amsmath}
\usepackage{graphicx}
\usepackage{booktabs}
\usepackage{array}
\usepackage{multirow}
\usepackage{xcolor}
\usepackage{url}

\title{AASSR: Knowledge-Parameterized Decision Making with Prophecy-Guided Imagination}

\author{
    \IEEEauthorblockN{Author Name}
    \IEEEauthorblockA{Affiliation\\
    Email}
}

\begin{document}
\maketitle

\begin{abstract}
This paper introduces AASSR/APASSR, a knowledge-parameterized decision-making
framework in which observations are not only stored as memory but reused as
parameters for future action templates. The framework represents discovered
information as KK/KV knowledge items, binds KV values into KK slots to generate
executable candidate actions, predicts candidate outcomes through a Prophecy
Module, and evaluates candidates before execution through an Imagination Cycle.
We evaluate the framework in controlled GridWorld environments that abstract
knowledge-action dependencies from the original pentesting motivation. In the
tested v2\_complex setting, the main C3 condition improves success rate and
reduces repeated actions compared with Random, PolicyABC, Q-learning, and a
partial-observation DQN baseline. Ablation studies suggest that repeat
suppression and error avoidance are especially important in the current
implementation.
\end{abstract}

\section{Introduction}

Many decision-making tasks cannot be reduced to selecting primitive actions.
In domains such as automated pentesting, actions are parameterized commands:
the agent must discover values such as target hosts, ports, services, and paths,
then reuse those values as parameters for future commands. This motivates a
decision-making framework in which knowledge is both memory and action material.

The central idea of AASSR/APASSR is:

\begin{quote}
Actions create knowledge, and knowledge creates the next actions.
\end{quote}

This paper evaluates that idea in GridWorld environments designed to abstract
knowledge-action dependencies.

\section{Framework}

\subsection{Knowledge Storage}

Knowledge Storage is represented as a typed parameter pool. KK denotes an
abstract knowledge key or action-template slot, and KV denotes a concrete value
that can be bound into that slot.

\subsection{Action Template Binding}

An action template such as \texttt{MOVE\_TOWARD \{KK\_FRONTIER\_CELL\}} becomes
an executable action when a concrete KV is selected from Knowledge Storage and
bound into the KK slot.

\subsection{PolicyABC}

PolicyABC decomposes action selection into WHAT, HOW, and WHERE axes.

\subsection{Prophecy Module}

The Prophecy Module predicts semantic knowledge change, error likelihood, and
goal relevance for candidate actions. Table, sequence, and Transformer-based
prophecy models are implementation variants rather than the framework itself.

\subsection{Imagination Cycle}

The Imagination Cycle evaluates candidate actions before execution using
Prophecy predictions. The current implementation performs depth-limited
candidate evaluation, not hidden-map simulation.

\section{Experimental Setup}

\subsection{Environments}

We evaluate on random\_key\_door, v2\_complex, and locked\_bottleneck
GridWorld environments.

\subsection{Conditions}

\begin{table}[h]
\centering
\caption{Experimental conditions.}
\begin{tabular}{ll}
\toprule
Condition & Description \\
\midrule
C0 & RandomScorer \\
C1 & PolicyABC \\
C2 & PolicyABC + Prophecy reward \\
C3 & PolicyABC + Prophecy Module + Imagination Cycle \\
C4 & Optional sequence-based Prophecy variant \\
QLEARN & Tabular Q-learning baseline \\
DQN\_PARTIAL & Partial-observation DQN baseline \\
ORACLE\_MDP & Full-map shortest-path upper bound \\
\bottomrule
\end{tabular}
\end{table}

\section{Results}

In the v2\_complex setting, C3 achieved higher success rate and lower repeat
rate than the Random, PolicyABC, Q-learning, and DQN\_PARTIAL baselines. The
ORACLE\_MDP condition is interpreted separately as a full-map upper bound.

\section{Ablation Studies}

We evaluate five ablation groups: Prophecy implementation, prediction-error
reward, Imagination depth/branching, Imagination mechanism terms, and Prophecy
score components.

\section{Discussion}

The results suggest that the current implementation benefits most from repeat
suppression and error avoidance rather than indiscriminate knowledge-gain
seeking. This supports the view that Knowledge Storage should be used to filter
and parameterize actionable candidates, not merely to maximize the amount of
newly discovered information.

\section{Limitations}

GridWorld is an abstract validation environment and does not directly prove
real-world pentesting performance. In addition, AASSR does not universally
dominate DQN in all environments or budgets.

\section{Conclusion}

AASSR/APASSR reframes knowledge storage as both memory and action-parameter
source. The framework links observation, knowledge update, action-template
binding, prediction, imagination, execution, and policy update into a closed
loop. Controlled GridWorld experiments support the usefulness of this structure
for knowledge-action dependency tasks.

\end{document}
